This paper develops a fixed-representation compression theory for zero-error identity recovery, governed by the collision fiber geometry $A_\pi$.

\paragraph{Compression contribution.}
The same fiber geometry determines the adaptive fiberwise budget, exact finite-block law, fiberwise decomposition of recoverable mass, and query-side counting law. The uniform single-block curve $D^\star(L)=\max(0,1-2^L/a)$ is the sharp closed-form specialization. If $\pi$ is injective, identity is free; if non-injective, the missing information must be paid in bits, queries, helper views, or distortion.

The query side ties directly back to the compression side. Any binary distinguishing family of size $d$ can realize at most $2^d$ transcripts on a collision fiber, so $\lceil \log_2 A_\pi \rceil \le d$ on the worst fiber. Complete derivability-aware query families reduce to orthogonal semantically minimal cores, on which exchange holds and matroid basis cardinality becomes the canonical query invariant.

The same ambiguity-fiber geometry governs when a representation is sufficient for exact downstream tasks, when helper views genuinely help, and when modular or factorized latents add exact information rather than architectural redundancy. In the learned-representation reading, injective representations support exact downstream identification with no extra metadata, while non-injective representations impose an irreducible side-information cost and, if that cost is left unpaid, a correspondingly irreducible deterministic distortion floor.

\paragraph{Systems payoff.}
Semantic abstraction imposes an exact identity cost: non-injective representations require auxiliary description with worst-case cost $\log_2 A_\pi$ bits and fiberwise cost $\lceil \log_2 |\pi^{-1}(u)| \rceil$. Symbolic handles are the standard payment mechanism. A neural encoder plus injective handle achieves zero-error recovery iff the handle distinguishes entities within each semantic fiber. \leanmeta{\LHrng{NSL}{1}{6}}

\paragraph{Secondary robustness consequence.}
Present collision-freedom need not persist under extension, and unrestricted future-generating code has no computable barrier-freedom certificate. Persistent zero-error guarantees require explicit growth models or identity metadata.

\paragraph{Mechanization.}
The main theorem chain is machine-checked in Lean 4, certifying the fiber-geometry framework, canonical query-core reduction, and neurosymbolic linking results in the finite explicit setting.

\subsection*{Acknowledgment: AI-use Disclosure}

Generative AI tools (including Codex, Claude Code, Augment, Kilo, and OpenCode) were used throughout this manuscript, across all sections (Abstract, Introduction, theoretical development, proof sketches, applications, conclusion, and appendix) and across all stages from initial drafting to final revision. The tools were used for boilerplate generation, prose and notation refinement, LaTeX and structure cleanup, translation of informal proof ideas into candidate formal artifacts (Lean and LaTeX), and repeated adversarial reviewer-style critique passes to identify blind spots and clarity gaps.

The author retained full intellectual and editorial control, including problem selection, theorem statements, assumptions, novelty framing, acceptance criteria, and final inclusion or exclusion decisions. No technical claim was accepted solely from AI output. Formal claims reported as machine-verified were admitted only after Lean verification (no \texttt{sorry} in cited modules) and direct author review; Lean was used as an integrity gate for responsible AI-assisted research. The author is solely responsible for all statements, citations, and conclusions.
